\section{The remaining R6 boundary: $a^2B$ infinitesimal}
\label{sec:r6-b-inf}

\paragraph{Source scope and relation to the preceding theorem.}
Retain the complete original weak-C2 source: $d=1$, $\lambda=\beta=0$,
$n\geq0$, the actual stationary contact, every strict gap, total scale,
and the complement of prior admissions. On R6, $t,h,l=1-h,r$ and
the original $j$ gap are positive bounded units, while $g,B,x=aB,y=tB$
are positive infinitesimals. This section treats $a$ unbounded and
$b=a^2B$ positive infinitesimal. The notation $s=1/a$, $w_s=ae$,
$B=s^2b$, $x=sb$, $e=s(1+hsb-g)$ is unchanged from (RB1).
The exact physical identity (RB2), equivalently RT2, remains in force.
It bounds $b$ on all unbounded-$a$ R6 and gives $\rho_0=1/2$ here,
where $\rho=n/e^2$. In particular this subdomain has $n>0$; the
literal $n=0$ source face is covered by the preceding $b$-unit theorem.

\begin{theorem}
On this entire original R6 source intersection with $b=a^2B$
infinitesimal, the complete original parent is a finite separated
source--testing tensor sum of squares at every independent projective
testing pair and every real or complex state. Both angle signs, all
literal projective endpoints, and all relative opening rates of
$s,b,g$ are included. No mixed inverse or mixed square root is used.
\end{theorem}

Together with the $b$-unit theorem and the inherited bounded-$a$
R6 theorem, this proves the entire R6 row. It does not by itself
admit another weak-C2 row or the global finite-lift theorem.

\paragraph{A regular source chart.}
Keep $R=1+hx$, $K=1+x$, $T=1+xr$, $k=\sqrt{x}$ and the contact
parameters $\delta=R/(K\sqrt{bt})$, $E=tK/R$ from (RB8). Define
\[
 \varepsilon=\delta^{-1},\qquad \jmath=\delta k,
 \qquad k=\varepsilon\jmath.
 \tag{RI1}
\]
Both $\varepsilon$ and $\jmath$ are positive infinitesimals. In this
section $\jmath$ denotes the contact scale, not the original source
gap called $j$. The exact source substitution is
\[
 \begin{gathered}
 x=\varepsilon^2\jmath^2,\quad K=1+x,\quad R=1+hx,\quad
 E=tK/R,\\
 b=\frac{\varepsilon^2R^2}{K^2t},\qquad
 s=\frac{t\jmath^2K^2}{R^2}.
 \end{gathered}                                                   \tag{RI2}
\]
All denominators in (RI2) are positive source units. There is no
comparison of a testing variable with either source infinitesimal.

Use precisely the complete contact-centered remainder $G$ in (RB11),
with the full coordinate inverse (RB7). Thus the two retained packets
on the contact chart are $(m-K_m)N^2$ and
$nD_uD_v(X+\gamma v\omega)^2$, where
\[
 p=(1+hx-g)/R,\quad c=l/(tT),\quad
 \gamma=rK/(sR),\quad D_u=1+u^2,\quad D_v=1+v^2.
\]
Its exact contact is now
\[
 (u,\omega,v;N,M_B)
   =(\varepsilon\jmath,\jmath,\varepsilon\jmath^2E;0,0).
 \tag{RI3}
\]
The six actual-contact checks in the preceding receipt prove the
vanishing of both normals and of the retained $n$ row here as well.

\paragraph{Two exact source cancellations.}
The quantities $mg$, $n\gamma$, $n\gamma^2$ and
$\alpha/(s^3b^2)$ have canceled denominators that are units at
$s=b=g=0$. Introduce
\[
 \mu_1=mg/R,\quad \mu_2=(mg)g/R^2,\quad
 F_0=p^2/2-n\gamma^2-2n\gamma p+\mu_2,\qquad
 A_N=p-2\mu_1-2n\gamma.
 \tag{RI4}
\]
Both $F_0$ and $A_N$ are divisible by $b$ in this rational unit ring.
Here is a short exact verification, independent of a series expansion.
At $b=0$, put $q=1-g$. Direct cancellation in the original source gives
\[
 \begin{gathered}
 p=q,\quad mg=\frac{q}{2(1+s)},\quad
 n\gamma=\frac{sq}{2(1+s)},\quad
 n\gamma^2=\frac{q(1-sq)}{2(1+s)}.
 \end{gathered}                                                    \tag{RI5}
\]
Substitution makes both expressions in (RI4) identically zero.
The denominators are units before this specialization; hence the
polynomial numerators are divisible by $b$, without introducing a
pole. Their needed residues are
\[
 \frac{F_0}{b}\sim\frac{h+t}{2},\qquad
 \left.\frac{A_N}{b}\right|_{s=g=0}=0,\qquad
 \left.\frac{A_N}{bs}\right|_{g=0}\sim h+t.
 \tag{RI6}
\]
The last quotient also has a unit canceled denominator at $s=b=0$.
Every statement in (RI5)--(RI6) is separately checked from the original
source formulas. In particular the source estimate is uniform in
$g/s$, $s/b$ and all other opening ratios.

\paragraph{Complete bounded raw support, including both first normals.}
Use the separated raw monomials
\[
 \mathsf A=\varepsilon\omega v,\quad
 \mathsf B=v,\quad \mathsf C=Eu\omega,\quad
 \mathsf D=\varepsilon^2\jmath^3E,
 \qquad \mathsf x=\mathsf B-\mathsf C.
 \tag{RI7}
\]
Set $H=\alpha+1/2$, $L=AD_v+Hv^2$, and $V_0=1-puv\omega$.
At zero normals, (RB7) gives
\[
 q_B=D_uvV_0,\quad uW-Y=cD_uvV_0,\quad
 Y=u\omega-cvV_0,\quad
 J=-(1-c)D_uvV_0,
\]
and the entire scalar has the following exact allocation:
\[
 G_0=f_A\mathsf A^2+f_{BB}\mathsf B^2+
 f_{BC}\mathsf B\mathsf C+f_{CC}\mathsf C^2+
 \varrho\mathsf D^2,                               \tag{RI8}
\]
where
\[
\begin{aligned}
 f_A={}&D_u^2D_vF_0/\varepsilon^2,\\
 f_{BB}={}&[Cc^2D_u^2D_v+(D+1/2)(1-c)^2D_u^2
                  +D_uLc^2]V_0^2\\
          &+2HD_ucV_0-D_u/2-D_u^2\mu_2v^2\omega^2,\\
 f_{BC}={}&[-2D_uLcV_0-2HD_u+2n\gamma D_uD_v]/E,\\
 f_{CC}={}&[D_uL+D_uD_vn\gamma^2v^2]/E^2,\\
 \varrho={}&\alpha D_u/\mathsf D^2.
\end{aligned}                                                      \tag{RI9}
\]
These coefficients are bounded. In particular
\[
 \frac{F_0}{\varepsilon^2}=
       \frac{R^2}{K^2t}\frac{F_0}{b},\qquad
 \frac{\alpha}{\mathsf D^2}=
       \frac1t\frac{\alpha}{s^3b^2}.
 \tag{RI10}
\]
The term that would give $u^2/\varepsilon^2$ in the first raw
coefficient has instead been placed in $f_{CC}$, using the exact
identity $u^2(v\omega)^2=(v\mathsf C/E)^2$.
The final term of $f_{BB}$ likewise allocates the whole
$-D_u^2\mu_2v^4\omega^2$ contribution without dividing by
$\varepsilon$. These are identities for the full scalar.

Write the exact normal expansion as
$G=G_0+L_NN+L_MM_B+G_{\rm hard}$, with $G_{\rm hard}$ quadratic.
The two complete linear coefficients are
\[
 L_N=R_A\mathsf A+R_B\mathsf B+R_C\mathsf C,
 \qquad L_M=S_B\mathsf B+S_C\mathsf C,               \tag{RI11}
\]
where
\[
\begin{aligned}
 R_A={}&D_uA_N/\varepsilon,\\
 R_B={}&D_u(p-2n\gamma)v^2\omega\\
 &-2uv[Cc^2D_uD_v+(D+1/2)(1-c)^2D_u+Ac^2D_v]V_0\\
 &-2Hcuv(1+cv^2V_0),\\
 R_C={}&2cuvL/E,\\
 S_B={}&2\{D_u[CcD_v-(D+1/2)(1-c)]+cL\}V_0+2H,\\
 S_C={}&-2L/E.
\end{aligned}                                                      \tag{RI12}
\]
Here $A_N/\varepsilon=
\varepsilon R^2(A_N/b)/(K^2t)$ is bounded. Thus no first-normal
pole has been discarded. Every hard coefficient is bounded as well:
the only testing inverse in the full change of state is $D_u^{-1}$,
and all original source coefficients in (RB11) are bounded after the
displayed cancellation. At $u=v=0$ the complete hard matrix is exactly
\[
 \operatorname{diag}(K_m+1/2,A+C+D+1/2).             \tag{RI13}
\]
The accompanying checker reconstructs $G$, differentiates in each
normal, and verifies (RI8)--(RI13) by multiplication back.

\paragraph{A bounded conductor for these complete coefficients.}
Define the contact differences
\[
 \mathsf y=\mathsf A-k\mathsf C,\quad
 \mathsf z=u\mathsf B-k\mathsf C,\quad
 \mathsf h=k\mathsf C-\mathsf D.
\]
The scaled conductor identities are
\[
\begin{aligned}
 \mathsf C(u-k)&=\mathsf z-u\mathsf x,\\
 \varepsilon\mathsf C(\omega-\jmath)
       &=\mathsf y-\varepsilon\omega\mathsf x,\\
 \mathsf C(v-\varepsilon\jmath^2E)
       &=E(\omega\mathsf z+\jmath\mathsf y-k\omega\mathsf x),\\
 \mathsf Cv(\omega-\jmath)
       &=E\omega(\omega-\jmath)\mathsf z+
         E\jmath\omega\mathsf y-Ek\omega^2\mathsf x.
\end{aligned}                                                       \tag{RI14}
\]
These are SE5--SE6 with $w=\omega$, also checked directly here.
For a monic contact division
\[
 S=(u-k)T_u+(\omega-\jmath)T_\omega+
           (v-\varepsilon\jmath^2E)T_v,
 \qquad T_\omega=\varepsilon U+vV,
\]
the exact bounded lift is
\[
\begin{aligned}
 \mathsf CS={}&\mathsf x[-uT_u-\varepsilon\omega U
                    -Ek\omega^2V-Ek\omega T_v]\\
 &+\mathsf y[U+E\jmath\omega V+E\jmath T_v]\\
 &+\mathsf z[T_u+E\omega(\omega-\jmath)V+E\omega T_v].
\end{aligned}                                                       \tag{RI15}
\]
No inverse of a small source parameter or a contact difference occurs.

We verify the hypothesis of (RI15) for the full coefficients, rather
than import the different small-$e$ source assumptions of SE8.
Clear the finite source unit denominators and $D_u^2$. They can be
chosen independent of $\omega$ and even in $u,v$. In (RI9), all
$\omega$ dependence is through $uv\omega$ and $v^2\omega^2$;
the other testing dependence is even in $u,v$. The same property
holds after subtracting any model of the form used below.

For the scalar square-lift algorithm CC4, substitute
$\mathsf A=\mathsf y+k\mathsf C$,
$\mathsf B=\mathsf x+\mathsf C$,
$\mathsf D=k\mathsf C-\mathsf h$ into (RI8).
The coefficient requiring its first conductor is
\[
 S=k^2f_A+f_{BB}+f_{BC}+f_{CC}+k^2\varrho.
\]
Its zero contact value follows by domain cancellation, as in CC4.
Divide monically in the order $u$, then $\omega$, then $v$.
After $u=k=\varepsilon\jmath$, the $\omega$ divided difference
belongs to $(\varepsilon v,v^2)$. Accordingly choose
$T_\omega=\varepsilon U+vV$ with $U,V\in(v)$.
Modulo $\varepsilon$, the remaining $v$ polynomial is even, so
$T_v\in(\varepsilon,v)$; also
$T_u|_{\varepsilon=u=v=0}=0$.

After this first conductor the coefficient of $\mathsf h$ remains
$-2k\varrho$. The sole gradient relation is
$(-2k,1,1,-1)$ on
$(\mathsf x,\mathsf y,\mathsf z,\mathsf h)$; the determinant of
the first three contact gradients is
$-\varepsilon^3\jmath^4E^2\ne0$. The adjugate cancellation of
CC5 therefore removes the contact values without putting an inverse
of this determinant into any coefficient. Subtract the contact
multiple $2k\varrho_c(-2k,1,1,-1)$.
Each of the four resulting coefficient remainders $J_i$ vanishes at
contact. Its restriction at $\varepsilon=u=v=0$ is independent of
$\omega$: for the $\mathsf x$ coefficient this follows from
$2f_{BB}+f_{BC}$; the other three restrictions are zero by the
$U,V,T_u,T_v$ properties just established and the factor $k$.
Consequently the next monic $\omega$ divided difference of every
$J_i$ belongs to $(\varepsilon,v)$. Formula (RI15) applies again.

The last scalar repair uses the exact circuit relation
\[
 k\mathsf C(\mathsf y+\mathsf z-\mathsf h-2k\mathsf x)
       =Q_1^2-\mathsf y\mathsf z,
\]
where $Q_1$ is the first of the separated radical pairings below.
This completes the same restricted scalar square lift as CC6.
For (RI11), one conductor suffices: after replacing
$\mathsf A,\mathsf B$ as above, the coefficient multiplying
$\mathsf C$, restricted at $\varepsilon=u=v=0$, is independent of
$\omega$ directly from (RI12). Its divided difference therefore
lies in $(\varepsilon,v)$ and (RI15) applies. The full scalar and
both full normals have thus been lifted over the bounded coefficient
ring. The restricted raw support $\mathsf A^2,\mathsf B^2,
\mathsf B\mathsf C,\mathsf C^2,\mathsf D^2$ is essential;
an unrestricted ideal claim is neither needed nor asserted.

The original stationary parent and the two retained contact-zero
packets give the exact scalar value and full first-jet identities
used here. The contact (RI3) is independent of $g$. Hence taking
the exact coefficient difference divisible by $g$ preserves these
identities and every raw-support condition.

\paragraph{Strict residues of the complete contact lift.}
Suppose now that $u,v,\omega$ are all positive infinitesimals, as
in the remaining contact state chart. Give
\[
 (\varepsilon,\jmath,u,\omega,v)
       \text{ weights }(1,1,2,1,3).
\]
Then $\mathsf x$ has weight three and the five circuit roots have
weight five. Every contact substitution and (RI14)--(RI15) preserves
the weights. Put $\kappa=(h+t)/(2t)$ and
\[
 \Lambda=\frac{1-s+s^2+u^2-[2l+h(h+t)]k^2}{2ht},\qquad
 \mu=\frac1{h(1+s)}.
\]
Direct exact collection of (RI9) through weight ten gives
\[
 G_0=\Lambda\mathsf x^2+\kappa T_4+O_{\rm wt}(12)
                   \quad\pmod g,                 \tag{RI16}
\]
where
\[
 T_4=\mathsf A^2+(u\mathsf B)^2+(k\mathsf C)^2+
          \mathsf D^2-4k^2\mathsf B\mathsf C.
\]
The two complete normal collections are
\[
 L_N=O_{\rm wt}(8),\qquad
 L_M=\mu\mathsf x+O_{\rm wt}(7)\quad\pmod g.
 \tag{RI17}
\]
For transparency, the only source terms needed for (RI16)--(RI17),
besides (RI6), are the following exact Taylor coefficients at $g=0$:
\[
\begin{aligned}
 A&=\frac t{2h}(1-s+s^2)+O_{\deg(s,b)}(3),\\
 C&=\frac{h+t}{2l}-\frac{t+1}{2l}s+
                      \frac{t+1}{2l}s^2+O_{\deg(s,b)}(3),\\
 D&=\frac{1-s+s^2}{2(t-l)}-
        \frac{sbt(h+t)}{2(t-l)^2}+O_{\deg(s,b)}(3),\\
 n\gamma&=\frac s2-\frac{s^2}2-\frac{sb(h+t)}2+
                         O_{\deg(s,b)}(3).
\end{aligned}                                                       \tag{RI18}
\]
Here the error denotes exact degree filtration in a rational ring
with unit denominators. Also $s=t\jmath^2+O_{\rm wt}(6)$,
$c=(l/t)(1-k^2)+O_{\rm wt}(6)$ and
$E=t(1+lk^2)+O_{\rm wt}(8)$. The symbolic checker verifies all
degree-two source derivatives, the complete scalar through weight
ten, and the complete $M_B$ normal through weight five. The weight
eight bound for $L_N$ follows term by term from (RI6), (RI12).
At $g=0$, after the displayed unit denominators are cleared, all
scalar weighted monomials have even weight, and all $M_B$ linear
monomials have odd weight. Hence the next possible degrees are
twelve and seven, respectively.

The stronger coefficient-wise statement also holds for the displayed
raw representation: $f_A-\kappa$ and $\varrho-\kappa$ have weight at
least two, while
\[
 f_{BB}-\Lambda-\kappa u^2,\qquad
 f_{BC}+2\Lambda+4\kappa k^2,\qquad
 f_{CC}-\Lambda-\kappa k^2
\]
each has weight at least six. The checker verifies their coefficients
through weight four separately. Similarly $S_B-\mu,S_C+\mu$ have
weight at least four separately, and $(R_A,R_B,R_C)$ have weights
at least $(3,5,5)$. Thus the chosen raw allocation of each exact
difference already has the claimed filtration; no low-weight
Gram coefficients are hidden by cancellations between raw products.

Apply the bounded conductor just proved to the exact differences in
(RI16)--(RI17). The scalar correction Gram coefficients have weight
at least two; the $M_B$ normal corrections have weight at least two;
the $N$ normal corrections have weight at least three. They are all
infinitesimal. The coefficients divisible by $g$ lift boundedly by the
same argument and are infinitesimal independently of every other
opening. Restoring the finite unit denominators preserves this fact.

For strictness use the two pairings on the separated monomials
\[
 a_i^2=(\mathsf A,u\mathsf B,k\mathsf C,\mathsf D),\qquad
 Q_1=a_1a_2-a_3a_4,\quad Q_2=a_1a_3-a_2a_4.
\]
Their product relation gives the exact two-pairing identity (RB15),
with the scaled roots. Its three-dimensional difference block has
determinant $1/2$, and both pairing pivots are positive. The complete
residue on
$(\mathsf x,\mathsf y,\mathsf z,\mathsf h,Q_1,Q_2,N,M_B)$
therefore consists of this strict circuit with coefficient
$\kappa>0$, the $N$ pivot $K_m+1/2>0$, and
\[
 \begin{pmatrix}
 1/(2ht)&1/(2h)\\
 1/(2h)&A_*+C_*+D_*+1/2
 \end{pmatrix},
 \quad A_*=t/(2h),\ C_*=(h+t)/(2l),\ D_*=1/[2(t-l)].
 \tag{RI19}
\]
Its determinant is exactly
\[
 \frac{C_*+D_*+1/2}{2ht}>0.                         \tag{RI20}
\]
All pivots are positive bounded units. The strict separated Gram
lemma WA10 yields a finite tensor-SOS identity for the complete
remainder. Restore the two actual nonnegative retained packets and
all separated clearing factors. This proves the entire local contact
chart, including every normal direction and arbitrary relative rates.

\paragraph{The $\omega$-unit state chart uses the original $n$ packet.}
When $u,v$ are infinitesimal, $Z=1$, and $\omega$ is a positive unit,
retain $nD_uD_vX^2$ itself, together with $(m-K_m)N^2$.
The complete remainder is (RB11) with its two negative $n\gamma$
and $n\gamma^2$ contributions absent. This choice is legitimate
directly from $n\geq0$ and does not require a contact lift away from
the contact state scale.

Use the raw roots
\[
 (A_o,\mathsf x,\mathsf D,N,M_B),\qquad A_o=\omega v.
\]
The only new inverse is $1/\omega$, which is pure testing and
bounded. The relations $\mathsf B=A_o/\omega$ and
$\mathsf C=A_o/\omega-\mathsf x$ put the entire scalar and both
first normals in this frame, by the same full formulas (RI9),
(RI12) with the absent $n$ terms deleted. All coefficients are
bounded; the constant term is still $\alpha D_u=\varrho\mathsf D^2$.
Explicitly $\mathsf A=\varepsilon A_o$, so the first scalar term
becomes $D_u^2D_v(p^2/2+\mu_2)A_o^2$ and the first $N$-normal
term becomes $D_u(p-2\mu_1)A_o$. The printed $\varepsilon$
denominators therefore cancel before any coefficient estimate.
Its exact residue is the $A_o$ pivot $1/2$, the $\mathsf D$ pivot
$\kappa$, the $N$ pivot $K_m+1/2$, and the coupled block (RI19).
The first $N$ normal has zero residue and the first $M_B$ normal
has residue $\mathsf x/h$. These complete-frame residues are
separate assertions in the receipt. Thus this entire state chart
also has a strict finite separated Gram.

\paragraph{All other independent tests and states.}
Every source bound used by the noncontact joins of the preceding
$b$-unit theorem is (RB4), which is valid here without a lower bound
on $b$. For bounded unit $u$ and infinitesimal $v$, retain the original
$n,\alpha$ events and use the exact full-state form (RB22), whose
determinant (RB23) is positive. For both angles infinitesimal, the
SD5 exact sign reduction disposes of the nonnegative-core branch
and literal zero axes by original-event squares. On the negative
branch choose the exact gauges $u,v>0$. If the normalized state has
$Z$ infinitesimal, BE4 gives the strict full-state frame
$(X,Y,uW,vW)$ using positive real caps of $m,A$.
Otherwise normalize $Z=1$; $\omega>0$ is bounded, and its two
possibilities are exactly the contact chart and the unit chart just
proved. No source/testing ratio split is used.

For finite remote $v_0\ne0$, use the actual centered-event identity
R5.5, the strict kernel determinant RT10 (printed as RB24), and a
sufficiently large positive real cap below unbounded $m$.
For unbounded $u$ and bounded unit $v$, use R6.5 specialized to
$n_0=0$ and the source residues (RB4). For unbounded $u$ and
infinitesimal $v$, use the exact BE $(\infty,0)$ full frame
$(p_uX+Z,W,Y,vX)$, $p_u=1/u$, with positive caps of $C,A$.
For unbounded $v$ and bounded $u$, R6.6 at $n_0=0$ is strict.
For both angles unbounded, R6.7 at $n_0=0$ has determinants
$D_*/2$ and $A_*(K_m+1/2)$ on its two nontrivial blocks.
These are the same exact homogeneous matrices listed in
\texttt{dependencies\_r6\_b\_unit.json}; their uses import the
identities, not a source theorem from another stratum.
Literal reciprocal zeros are polynomial values of those charts.
All signs and complex states follow by the same exact gauges and
the real/imaginary split. This exhausts all independent tests.

\paragraph{Additional exact dependencies and audit scope.}
Beyond the preceding dependency index, the scaled algebra is SE5--SE7
in \texttt{inherited\_r18/inherited\_r17/inherited\_r16/
inherited\_r13/small\_e\_all\_r\_contact.tex} relative to the R19
root. This section proves its required rational-support hypotheses
afresh in (RI8)--(RI15); the small-$e$ native source theorem, its
physical parameter constraints, and its particular residues are not
imported. CC4--CC6 supplies the monic square-lift mechanism and the
gradient/repair argument. The checker corroborates exact algebra;
the bounded polynomial support proof and exhaustive chart join are
the mathematical argument above. Neither the checker nor this local
R6 result claims a global finite lift or external peer review.
The complete exact path list, including this additional algebraic
dependency, is in \texttt{dependencies\_r6\_b\_inf.json}.
